\documentclass{article}
\usepackage{amsmath, amssymb}

\title{1.2.6-10}

\begin{document}
\maketitle

\subsubsection*{a)}

$n = pq + r$

$\binom{n}{p} = \frac{(pq+r)(pq+r-1)\cdots(pq+1)pq(pq-1)\cdots(pq-(p-r-1)}{p!} = \langle\frac{q(pq+r)(pq+r-1)\cdots(pq+1)(pq-1)\cdots(pq-(p-r-1)}{(p-1)!}\rangle_p =$

$= \langle q \rangle_p \cdot \langle(p-1)!\rangle_p \cdot \langle((p-1)!)^{-1}\rangle_p = \langle q \rangle_p = \langle \lfloor\frac{n}{p}\rfloor \rangle_p$

\subsubsection*{b)}

$\binom{p}{k} = \frac{p(p-1)(p-2)\cdots(p-k+1)}{k!} = p\langle(p-1)(p-2)\cdots(p-k+1)\rangle_p \cdot \langle(k!)^{-1}\rangle_p = \langle 0 \rangle_p$

\subsubsection*{c)}

$\binom{p-1}{k} = \frac{(p-1)(p-2)\cdots(p-k)}{k!} = \langle(-1)^k k!\rangle_p \cdot \langle(k!)^{-1}\rangle_p = \langle(-1)^k \rangle_p$

\subsubsection*{d)}

$\binom{p+1}{k} = \frac{(p+1)p(p-1)\cdots(p-k+2)}{k!} = p\langle(-1)^{k-2}(k-2)!\rangle_p \cdot \langle(k!)^{-1}\rangle_p = \langle 0 \rangle_p$

\subsubsection*{e)}

$n = pq + r$

$k = ps + t$

$k! = (ps+t)! =$

$= 1 \cdot 2 \cdot \ldots \cdot (p-1) \cdot p \cdot$

$\cdot (p+1) \cdot (p+2) \cdot \ldots \cdot (p + p-1) \cdot 2p \cdot$

$\cdots$

$\cdot (p(s-1)+1) \cdot (p(s-1)+2) \cdot \ldots \cdot (p(s-1) + p-1) \cdot ps \cdot$

$\cdot (ps+1) \cdot (ps+2) \cdot \ldots \cdot (ps + t) =$

$= (\langle(p-1)!\rangle_p)^s \cdot p^s s! \cdot \langle t!\rangle_p = (\langle-1\rangle_p)^s \cdot p^s s! \cdot \langle t!\rangle_p$

Similarly, $n! = (\langle-1\rangle_p)^q \cdot p^q q! \cdot \langle r!\rangle_p$

$n - k = p(q-s) + (r - t)$

\paragraph*{1. \$\{\} r - t \textbackslash\{\}geq 0 \{\}\$}

$(n-k)! = (\langle-1\rangle_p)^{q-s} \cdot p^{q-s} (q-s)! \cdot \langle (r-t)!\rangle_p$

$\binom{n}{k} = \frac{n!}{k!(n-k)!} = \frac{(\langle-1\rangle_p)^q \cdot p^q q! \cdot \langle r!\rangle_p}{(\langle-1\rangle_p)^s \cdot p^s s! \cdot \langle t!\rangle_p \cdot (\langle-1\rangle_p)^{q-s} \cdot p^{q-s} (q-s)! \cdot \langle (r-t)!\rangle_p} = \langle\binom{q}{s}\binom{r}{t}\rangle_p = \langle\binom{\lfloor\frac{n}{p}\rfloor}{\lfloor\frac{k}{p}\rfloor}\binom{n \mod p}{k \mod p}\rangle_p$

\paragraph*{2. \$\{\} r - t < 0 \{\}\$}

$\binom{r}{t} = 0$

$n - k = p(q-s-1) + (p + r -t)$

$\binom{n}{k} = \frac{n!}{k!(n-k)!} = \frac{(\langle-1\rangle_p)^q \cdot p^q q! \cdot \langle r!\rangle_p}{(\langle-1\rangle_p)^s \cdot p^s s! \cdot \langle t!\rangle_p \cdot (\langle-1\rangle_p)^{q-s-1} \cdot p^{q-s-1} (q-s-1)! \cdot \langle (p+r-t)!\rangle_p} = \langle 0 \rangle_p = \langle\binom{\lfloor\frac{n}{p}\rfloor}{\lfloor\frac{k}{p}\rfloor}\binom{n \mod p}{k \mod p}\rangle_p$

\subsubsection*{f)}

$\binom{n}{k} \overset{e)}{\equiv} \binom{a_rp^{r-1} + a_{r-1} p^{r-2} + \cdots + a_1}{b_rp^{r-1} + b_{r-1} p^{r-2} + \cdots + b_1} \cdot \binom{a_0}{b_0} \equiv \binom{a_rp^{r-2} + a_{r-1} p^{r-3} + \cdots + a_2}{b_rp^{r-2} + b_{r-1} p^{r-3} + \cdots + b_2} \cdot \binom{a_1}{b_1} \cdot \binom{a_0}{b_0} \equiv \cdots \equiv \binom{a_r}{b_r} \cdots \binom{a_0}{b_0}\ (\mathrm{modulo}\ p)$


\end{document}
